\documentclass[10pt,a4paper]{amsart}
\usepackage[margin=19mm]{geometry}
\usepackage{amsmath,amssymb,mathtools}
\usepackage[scr=rsfs,cal=euler]{mathalfa}
\usepackage{xfrac}
\usepackage[unicode,hidelinks,bookmarks=false]{hyperref}
\newcommand{\ie}{i.e.\ }
\newcommand{\cf}{cf.\ }
\newcommand{\resp}{respectively\ }
\newcommand{\Subsection}{\subsection}
\providecommand{\nobreakdash}{\nobreak\hbox{-}}
\setlength{\emergencystretch}{2em}
\multlinegap0pt
\usepackage{bm}
\usepackage{bbm}
\usepackage{dsfont}
\usepackage{xspace}
\usepackage{cancel}
\usepackage{mathtools}
\usepackage{amsthm}
\makeatletter
\@ifundefined{thm}{\newtheorem{thm}{Thm}}{}
\@ifundefined{assumption}{\newtheorem{assumption}{Assumption}}{}
\@ifundefined{lemma}{\newtheorem{lemma}[thm]{Lemma}}{}
\@ifundefined{proposition}{\newtheorem{proposition}[thm]{Proposition}}{}
\makeatother
\newcommand{\RR}{\mathbb{R}}
\newcommand{\cG}{{\mathcal{G}}}
\newcommand{\norm}[1]{\| #1 \|}
\title[On the Linear Speedup of the Push-Pull Method for Decentrali]{On the Linear Speedup of the Push-Pull Method for Decentralized Optimization over Digraphs}
\author{Liyuan Liang, Gan Luo, Kun Yuan}
\date{Source: arXiv:2506.18075v5 (CC BY 4.0)}
\begin{document}
\maketitle

\noindent\emph{Source and licence.} On the Linear Speedup of the Push-Pull Method for Decentralized Optimization over Digraphs. Version arXiv:2506.18075v5 (CC BY 4.0), distributed under the Creative Commons Attribution 4.0 International licence, as recorded in the arXiv licence metadata for this exact version and shown on the abstract page https://arxiv.org/abs/2506.18075v5. Full rights notice: licenses/arxiv-2506.18075-CC-BY-4.0.txt.\par\smallskip

\noindent\textit{Editorial scope.} Counted unit: the Lemma whose source label is \texttt{lem:rolling sum} in the article (source0.tex lines 903--913, source label \texttt{lem:rolling sum}), delivered together with the article's own complete original proof of it (source0.tex lines 914--964, one \texttt{proof} environment of 51 lines). No mathematics is written, repaired or re-derived: statement and proof are the article's own text line for line. Required context retained uncounted: ass:graph (lines 342--344); ass:matrix (lines 346--348); prop:perron (lines 352--362); eq:s\_A definition (lines 400--404); prop:s\_A,s\_B (lines 407--415). Every definition, display and label of those lines is retained unchanged. Omitted from this package and not redistributed: 1--341, 345--345, 349--351, 363--399, 405--406, 416--902, 965--1539. Nothing inside a retained excerpt is omitted or altered: every retained body is delivered exactly as the article's lines are written, so no omission receipt of any kind accompanies this package. Citations: the retained excerpts carry no citation command at all, so no bibliography entry of the article is transcribed and no reference is redistributed. Reference adaptation: none was needed and none was made; every reference inside the delivered unit resolves inside it, so no unresolved reference or citation is produced. Printed slips kept literal: no printed word, symbol, hypothesis or number of the counted statement or of its proof was altered. Layout and class: the article's own class is \texttt{article} with packages including geometry, hyperref, url, amsmath, amssymb, amsthm, amscd, mathrsfs, amsfonts, inputenc, graphicx, booktabs, indentfirst, enumerate; the wrapper is the lane's pinned \texttt{amsart} editorial wrapper at 10pt on A4, which restates the theorem-like environments the retained text needs and adds the article's own macro definitions that the retained text needs (\texttt{\textbackslash RR}, \texttt{\textbackslash cG}, \texttt{\textbackslash norm}), with extra wrapper packages (bm, bbm, dsfont, xspace, cancel, mathtools). The delivered numbering is the unit's own. Assets: no retained span contains a figure environment, a graphics-inclusion command, a PDF-page inclusion, a table or a TikZ picture; no image file is copied into this package, the package declares no asset dependency and no image file is redistributed with it. Licence: version arXiv:2506.18075v5 of this article is distributed under the Creative Commons Attribution 4.0 International licence, as recorded in the arXiv licence metadata for this exact version and shown on the abstract page https://arxiv.org/abs/2506.18075v5. A full rights notice is delivered at licenses/arxiv-2506.18075-CC-BY-4.0.txt. Characters: every retained excerpt is pure ASCII, so the delivered engine, XeLaTeX with its default Latin Modern text font, reports no missing character.
\begin{assumption}[\sc \small Communication Graph]\label{ass:graph}
     Assume $\cG_A$ and $\cG_B$ to be strongly connected digraphs (see the definition in notations) with $n$ nodes.
\end{assumption}

\begin{assumption}[\sc \small Weight Matrix]\label{ass:matrix}
    The weight matrix \(A\) is row-stochastic with underlying digraph \(\mathcal{G}_A\), and the weight matrix \(B\) is column-stochastic with underlying digraph \(\mathcal{G}_B\). Both \(A\) and \(B\) have positive trace.
\end{assumption}

\begin{proposition}[\sc \small Exponential Decay]\label{prop:perron}
For \(A\) and \(B\) satisfying Assumptions~\ref{ass:graph} and~\ref{ass:matrix}, there exist entrywise positive Perron vectors \(\pi_A, \pi_B \in \mathbb{R}^n\) such that
\[
\pi_A^\top A = \pi_A^\top, \quad B \pi_B = \pi_B, \quad \norm{\pi_A}_1=\norm{\pi_B}_1= 1, \quad \pi_A^\top\pi_B >0.
\]
Besides, there exist constants $p_A, p_B>0$, $\lambda_A,\lambda_B\in [0,1)$ such that 
\begin{equation}\label{eq:exp decay}
    \norm{A^t-A_\infty}_2\le p_A\lambda_A^t,\quad\norm{B^t-B_\infty}_2\le p_B\lambda_B^t,\quad \forall k\ge 0.
\end{equation}

\end{proposition}

\begin{align}
    s_A &:= \max\big\{\sum_{i=0}^\infty \norm{A^i - A_\infty}_2, \sum_{i=0}^\infty \norm{A^i - A_\infty}_2^2\big\},\label{eq:s_A definition}\\
    s_B &:= \max\big\{\sum_{i=0}^\infty \norm{B^i - B_\infty}_2,\sum_{i=0}^\infty \norm{B^i - B_\infty}_2^2\big\},\label{eq:s_B definition} \vspace{2mm}\\ 
    c &:= n\, \pi_A^\top \pi_B.\label{eq:c definition}
\end{align}

\begin{proposition}[\sc\small Relation with existing metrics]\label{prop:s_A,s_B}
Suppose Assumption~\ref{ass:graph} and \ref{ass:matrix} hold.
The following inequalities on $s_A$, $s_B$ and $c$ hold:
    \begin{align*}
    1\le s_A\le \frac{M_A^2(1+\frac{1}{2}\ln(\kappa_A))}{1-\beta_A},&\quad 1\le s_B \le \frac{M_B^2(1+\frac{1}{2}\ln(\kappa_B))}{1-\beta_B},\\
    n\max\{\min \pi_A, \min \pi_B\}&\le  \,c \le n\min\{\max \pi_A, \max \pi_B\}. 
\end{align*}
where \( M_A := \max_{t \ge 0} \|A^t - A_\infty\|_2 \) and \( M_B := \max_{t \ge 0} \|B^t - B_\infty\|_2 \) are finite numbers which can be proved to be no larger than $\sqrt{n}$. 
\end{proposition}

\begin{lemma}[\sc \small Rolling Sum Lemma]\label{lem:rolling sum}
For a row-stochastic weight matrix satisfying Proposition~\ref{prop:perron}, the following inequality holds for any $ T\ge 0$:
    \begin{align}\label{eq:rolling sum A}
    \sum_{k=0}^T\norm{\sum_{i=0}^k(A^{k-i}-A_\infty)\Delta^{(i)}}_F^2\leq s_A^2\sum_{i=0}^T\norm{\Delta^{(i)}}_F^2,
    \end{align}
    where $\Delta^{(i)}\in \RR^{n\times d}$ are arbitrary matrices, $s_A$ is defined as in~\eqref{eq:s_A definition}. Furthermore, if Assumption~\ref{ass:graph} and~\ref{ass:matrix} hold, we have
    \begin{align*}
        s_A:= \max\{\sum_{i= 0}^\infty\norm{A^i-A_\infty}_2, \sum_{i= 0}^\infty\norm{A^i-A_\infty}_2^2\}\le \frac{M_A^2(1+\frac{1}{2}\ln(\kappa_A))}{1-\beta_A}.
    \end{align*} 
     where $\beta_A$, $\kappa_A, M_A$ are defined in Proposition~\ref{prop:s_A,s_B}. The row-stochastic matrix \(A\) and the subscript $A$ can be replaced by a column-stochastic matrix \(B\) without changing the lemma.
\end{lemma}

\begin{proof}
First, we prove that
\begin{align}\label{eq: rolling sum lemma-1}
    \norm{A^i-A_\infty}_2\le \sqrt{\kappa_A}\beta_A^{i}, \forall i \ge 0.
 \end{align}
Notice that $\beta_A:=\norm{A-A_\infty}_{\pi_A}$ and \[\norm{A^i-A_\infty}_{\pi_A}= \norm{(A-A_\infty)^i}_{\pi_A}\le \norm{A-A_\infty}^i_{\pi_A}=\beta_A^i,\]
we have
\[
\norm{(A^{i}-A_\infty)v}=\norm{\Pi_A^{-1/2}(A^i-A_\infty)v}_{\pi_A}\le \sqrt{\underline{\pi_A}}\beta_A^i\norm{v}_{\pi_A}\le \sqrt{\kappa_A}\beta_A^i\norm{v},
\]
where $\underline{\pi_A}:=\min \pi_A$, $\overline{\pi_A}:=\max \pi_A$.
The last inequality comes from $\norm{v}_{\pi_A}\le \overline{\pi_A} \norm{v}$. Therefore, \eqref{eq: rolling sum lemma-1} holds.

Second, we want to prove that
\begin{align}\label{eq: rolling sum lemma-2}
    s_A\le \frac{M_A^2(1+\frac{1}{2}\ln(\kappa_A))}{1-\beta_A}.
\end{align}
Towards this end, we define $M_A:=\max\{1, \max_{k\ge 0}\norm{A^k-A_\infty}_2 \}$. According to \eqref{eq: rolling sum lemma-1}, $M_A$ is well-defined. We also define $p=\max\left\{\frac{\ln(\sqrt{\kappa_A})-\ln(M_A)}{-\ln(\beta_A)}, 0\right\}$,
then we can verify that $\norm{A^i-A_\infty}_2\le \min\{M_A, M_A\beta_A^{i-p}\},\forall i\ge 0.$ With this inequality, we can bound $\sum_{i=0}^k\norm{A^{k+1-i}-A_\infty}_2$ as follows:
\begin{align}\label{eq: rolling sum lemma-3}
    &\sum_{i=0}^k\norm{A^{k-i}-A_\infty}_2=\sum_{i=0}^{\min\{\lfloor p \rfloor, k\}}\norm{A^i-A_\infty}_2+ \sum_{i=\min\{\lfloor p \rfloor, k\}+1}^{k}\norm{A^{i}-A_\infty}_2\\
    \le &\sum_{i=0}^{\min\{\lfloor p \rfloor, k\}} M_A+\sum_{i=\min\{\lfloor p \rfloor, k\}+1}^{k} M_A\beta_A^{i-p}\nonumber\\
    \le & M_A\cdot (1+\min\{\lfloor p \rfloor, k\})+M_A\cdot \frac{1}{1-\beta_A}\beta_A^{\min\{\lfloor p \rfloor, k\}+1-p}.\nonumber
\end{align}
If $p=0$, \eqref{eq: rolling sum lemma-3} is simplified to $\sum_{i=0}^k\norm{A^{k-i}-A_\infty}_2\le M_A\cdot\frac{1}{1-\beta_A}$ and \eqref{eq: rolling sum lemma-2} is naturally satisfied.
If $p>0$, let $x=\min\{\lfloor p \rfloor, k\}+1-p\in [0,1)$, \eqref{eq: rolling sum lemma-2} is simplified to
\[
\sum_{i=0}^k\norm{A^{k-i}-A_\infty}_2\le M_A\!\left(x+p+\frac{\beta_A^x}{1-\beta_A}\right)\le M_A\!\left(p+\frac{1}{1-\beta_A}\right)=\frac{M_A(1+\frac12 \ln(\kappa_A))}{1-\beta_A}.
\]
By the definition of $p$, we have $p\le \frac{\frac{1}{2}\ln(\kappa_A)}{1-\beta_A}$. This completes the estimation of $\sum_{i=0}^{\infty}\norm{A^{k-i}-A_\infty}_2$. For the summation of the squared norm, we can use the same $p$ and obtain the upper bound without modifying the previous steps:
\[
\sum_{i=0}^k\norm{A^{k-i}-A_\infty}_2^2\le M_A^2\!\left(p+\frac{1}{1-\beta_A}\right)=\frac{M_A^2(1+\frac12 \ln(\kappa_A))}{1-\beta_A}.
\]
This finishes the proof of \eqref{eq: rolling sum lemma-2}.

Finally, to obtain \eqref{eq:rolling sum A}, we use Jensen's inequality. For any positive numbers $a_0,\ldots,a_k$ satisfying $\sum_{i=0}^{k} a_i=1$, we have
\begin{align}\label{eq: rolling sum lemma-4}
     &\norm{\sum_{i=0}^k(A^{k-i}-A_\infty)\Delta^{(i)}}_F^2=\norm{\sum_{i=0}^k a_{k-i}\cdot a_{k-i}^{-1}(A^{k-i}-A_\infty)\Delta^{(i)}}_F^2\\
    \le& \sum_{i=0}^k a_{k-i}\norm{a_{k-i}^{-1}(A^{k-i}-A_\infty)\Delta^{(i)}}_F^2\le \sum_{i=0}^k a_{k-i}^{-1}\norm{A^{k-i}-A_\infty}_2^2\norm{\Delta^{(i)}}_F^2.\nonumber
\end{align}
By choosing $a_{k-i}=(\sum_{i=0}^k\norm{A^{k-i}-A_\infty}_2)^{-1}\norm{A^{k-i}-A_\infty}_2$ in \eqref{eq: rolling sum lemma-4}, we obtain that
\begin{align}\label{eq: rolling sum lemma-5}
    \norm{\sum_{i=0}^k(A^{k-i}-A_\infty)\Delta^{(i)}}_F^2\le \sum_{i=0}^k\norm{A^{k-i}-A_\infty}_2\cdot\sum_{i=0}^k\norm{A^{k-i}-A_\infty}_2\norm{\Delta^{(i)}}_F^2.
\end{align}
By summing up \eqref{eq: rolling sum lemma-5} from $k=0$ to $T$, we obtain that
\begin{align*}
    &\sum_{k=0}^T\norm{\sum_{i=0}^k(A^{k-i}-A_\infty)\Delta^{(i)}}_F^2 \le s_A\sum_{k=0}^T\sum_{i=0}^k\norm{A^{k-i}-A_\infty}_2\norm{\Delta^{(i)}}_F^2\\
    \le& s_A \sum_{i=0}^T\!\left(\sum_{k=i}^T\norm{A^{k-i}-A_\infty}_2\right)\norm{\Delta^{(i)}}_F^2\le s_A^2\sum_{i=0}^T\norm{\Delta^{(i)}}_F^2,
\end{align*}
which finishes the proof of this lemma.
\end{proof}

\end{document}
